\documentclass[a4paper]{article}
\usepackage[pdftex]{graphicx}
\usepackage[utf8]{inputenc}
\usepackage{enumerate}
\usepackage{siunitx}
\sisetup{locale=DE} 
\usepackage{href-ul}
\hypersetup{
	colorlinks=true,
	linkcolor=blue,
	urlcolor=blue}
\usepackage{icomma}
\usepackage{amssymb}
\usepackage{geometry}
\geometry{a4paper, top=15mm, left=15mm, right=15mm, bottom=15mm,
	headsep=10mm, footskip=12mm}

\begin{document}
{\bf Kräftezerlegung}
\begin{enumerate}[1.]


{\bf \item Auf einer Rampe mit der Länge von $l=\SI{3,0}{\meter}$ wird ein Fass mit der Gewichtskraft $F_G =\SI{600}{\newton}$ auf eine $h =\SI{1,20}{\meter}$ hohe Bühne gerollt.} Ermittle mit einem Kräfteplan die erforderliche Zugkraft $ F_Z $.

\begin{minipage}{0.6\textwidth}
{\bf \item Ein LKW der Masse $m=\SI{28}{\tonne}$ fährt eine Passstraße mit der Steigung  10\% bergauf.} Nach einer Fahrstrecke von 20 km hat er die Passhöhe erreicht. \\
\end{minipage}
\hfill
\begin{minipage}{0.4\textwidth}
\includegraphics[width=6 cm]{schieb115.pdf}
\end{minipage}

\begin{enumerate}[a)]
\item Ermittle näherungsweise den Höhenunterschied, den der LKW bei dieser Fahrt zurücklegt.
\item Berechne die Kraft, mit der der LKW bei dieser Bergfahrt durch den Motor nach oben gezogen wird (falls keine Reibung vorhanden wäre).
\end{enumerate}

\begin{minipage}{0.5\textwidth}
{\bf \item An einer Hauswand soll, wie nebenstehend in der maßstabsgerechten Zeichnung dargestellt, ein Schild angebracht werden:}
Um zu vermeiden, dass sich das Schild durch sein eigenes Gewicht ($F_G=\SI{52}{\newton}$) senkt, soll zwischen den Punkten A und B ein Seil gespannt werden. Die waagrechte Stange, die das Schild trägt, ist $l=\SI{2,0}{\meter}$ lang ; ihre Masse wird vernachlässigt. An der Mauer ist sie mit einem Scharnier befestigt. Der Schwerpunkt S des Schildes  hat von der Mauer den Abstand $a=\SI{1,00}{\meter}$. Die obere Halterung des Seils liegt $h=\SI{1,50}{\meter}$ oberhalb des mauerseitigen Endes der Stange.
\end{minipage}
\hfill
\begin{minipage}{0.45\textwidth}
\includegraphics[width=5.4 cm]{schild114.pdf}
\end{minipage}

Ergänze einen Kräfteplan in der Zeichnung und ermittle, welche Kraft das Seil aufnehmen muss.
$[\SI{1}{\centi\meter} \ \textnormal{Pfeillänge} \ \hat{=}\ \SI{10}{\newton}]$ 


\begin{minipage}{0.65\textwidth}
{\bf \item  Schiefe Ebene; Kraftzerlegung}\\
Viola steht mit ihrer kompletten Skiausrüstung (Gesamtmasse $m=\SI{51}{\kilogram}$) oben am Berg und stürzt sich mutig einen steilen Hang hinab. Der Neigungswinkel $\alpha$ beträgt $37^{\circ}$.
\end{minipage}
\hfill
\begin{minipage}{0.65\textwidth}
\includegraphics[width=4 cm]{violaski129.pdf}
\end{minipage}

\begin{enumerate}[a)]
\item Erstelle einen Kräfteplan mit der Gewichtskraft  $\vec{F_G}$, mit
 der Kraft $\vec{F_H}$, welche sie in Richtung Tal beschleunigt und mit der Kraft  $\vec{F_N}$, die auf die Schneedecke ausgeübt wird.\\ $[\vec{g}=\SI{9,81}{\meter\over {\second^2}}]$ 
\item Wie groß ist die Beschleunigung $\vec{a}_0$, die Viola bergab erfährt, wenn Reibungskräfte unberücksichtigt bleiben?
\item Nun sei die Gleitreibungszahl $ \mu = 0,15 $. Wie groß ist nun die Beschleunigung
 $\vec{a}_R$?   
\end{enumerate}

\end{enumerate} 

\fbox{
	\begin{minipage}{0.5\textwidth}
		Zur Lösung bitte \href{https://www.okuyakl.de/phys/p9kraL404/ll404.pdf}{hier klicken} oder den QR-Code scannen.\\
	Weitere Arbeitsblätter gibt es unter 
	
	\href{https://www.okuyakl.de}{www.okuyakl.de}
	\end{minipage}
	\hfill
	\begin{minipage}{0.4\textwidth}
		\includegraphics[width=1.5 cm]{../../viecher/zwe03}
		\includegraphics[width=3 cm]{qr404}
		\includegraphics[width=2 cm]{../../viecher/afanticon1}
		
	\end{minipage}}

\end{document}%Lösung-------------------------------------------
